

\begin{frame}
    \titlepage
\end{frame}

\begin{frame}{Introduction}
    \begin{itemize}
        \item Rock-Paper-Scissors is not random.
        \item Humans are not random, so given a series of games of
        rock-paper-scissors, you can predict what moves will be played.
        % \item We do that.
    \end{itemize}
\end{frame}

\begin{frame}{Data Collection}
    \begin{itemize}
        \item There's no data we could find on rock-paper-scissors games
        \item We have a program to collect our own data.

        \item The user plays a series of games against a random opponent.
        \item We collect the games and store them for later processing.
    \end{itemize}
\end{frame}

\begin{frame}{Data Format}
    \begin{itemize}
        \item We store games as 2d arrays, with each play stored as a 3-element array.
        \item We store series of games from the same session in lists.
    \end{itemize}
\end{frame}

\begin{frame}{Data Format}
    \[
    \begin{bmatrix}
        \begin{bmatrix}
            \begin{bmatrix}\tikzmarknode{game1a}{1} & \tikzmarknode{game1b}{0} & \tikzmarknode{game1c}{0}\end{bmatrix}, & \begin{bmatrix}\tikzmarknode{game1d}{0} & \tikzmarknode{game1e}{1} & \tikzmarknode{game1f}{0}\end{bmatrix}
        \end{bmatrix}, \\\\
        \begin{bmatrix}
            \begin{bmatrix}\tikzmarknode{game2a}{0} & \tikzmarknode{game2b}{0} & \tikzmarknode{game2c}{1}\end{bmatrix}, & \begin{bmatrix}\tikzmarknode{game2d}{1} & \tikzmarknode{game2e}{0} & \tikzmarknode{game2f}{0}\end{bmatrix}
        \end{bmatrix}
    \end{bmatrix}
    \]

    \only<2>{
        \begin{tikzpicture}[overlay, remember picture]
            \draw[<-, thick, red] (game1a.north) -- ++(0, 0.5) node[above] {Rock};
            \draw[<-, thick, red] (game1e.north) -- ++(0, 0.5) node[above] {Paper};
            \draw[<-, thick, blue] (game2c.south) -- ++(-0.5, -0.5) node[below] {Scissors};
            \draw[<-, thick, blue] (game2d.south) -- ++(0.5, -0.5) node[below] {Rock};
            \draw[<-, thick, black] (game1f.east) -- ++(1, 0) node[right] {Game 1};
            \draw[<-, thick, black] (game2f.east) -- ++(1, 0) node[right] {Game 2};
        \end{tikzpicture}
    }
\end{frame}

\begin{frame}{Methodology}
\begin{itemize}
    \item Classifying methodology: Computer play weights are initialized as
    pure random (.33, .33, .33).
    \item Apriori runs, then the most frequent item is found, where item is
    a single type of game (e.g. human plays rock, computer plays paper).
    \item Then, depending on whether that most frequent item is a winning or
    losing game, we update the weights of the chances for the computer to
    play an item in proportion to the frequency of the specific game.
\end{itemize}
\end{frame}



\begin{frame}{Preliminary Results}
 % \begin{block}{Findings}
 % Initial analysis shows significant bias in opening moves.
    As we have not fully implemented our methodology, current results show
    little to no improvement over random, however as we continue to implement our design,
    this is sure to change.
 % \end{block}
\end{frame}

\begin{frame}{Conclusion}
    Any questions?
\end{frame}
